% !TeX program = xelatex
% ============================================================================
%  第 14 课　复习、总结与展望　—— 教师版讲义
%  与 02-课件/第14课-复习总结与展望/第14课-复习总结与展望.tex 同步
% ============================================================================

\jhlesson{14}{复习、总结与展望}{时长 45 分钟\quad·\quad 教具：无\quad·\quad 本课不发单页}

\begin{mubiao}
  \begin{itemize}
    \item 能自己说出这门课学过的三类「绕圈」。
    \item 能指着主线图讲出它和群的四条规则的关系。
    \item 知道群之后还有环和域，知道五次方程的故事。
  \end{itemize}
\end{mubiao}

\noindent{\small 本课节奏：0—5 开场｜5—26 板书主线图｜26—36 抽查｜36—45 展望（环与域 2 分钟｜四次方程的公式 3 页 ≤ 4 分钟｜五次方程的故事 2 分钟｜最后一行 1 分钟）。}

\begin{beifen}
  \textbf{整门课的最后一道保险}：这门课能排满 14 课时已属不易——最后三周要留给期末复习，校本课有可能被停，
  **所以费马小定理必须在第 12 课讲完**，第 13、14 课是"能上就上、上不了不欠账"。
  如果这一课真的上不成，**第 12 课的收尾（定理原文 + 黑板右边那行字）就是全课的自然终点**，
  不必另外补课。
\end{beifen}

\begin{beifen}
  \textbf{时间不够}（这节课最容易超时的是"展望"那 9 分钟里塞了四件事）：
  砍掉三次、四次方程那两页，只讲二次那一页（一句话带过"三次四次也有，长这样，课后看"）；
  抽查从 8—10 个点减到 4—5 个，只点每条线一个点。
  \textbf{最后一行板书「群不是新东西，它就是『绕圈』这件事的名字」不许砍。}
  \textbf{时间富余}：方程那三页多讲两句（"公式长成这样，数学家花了多少年"）；
  或者让学生说一说"这门课里哪一件事最让你意外"。
\end{beifen}

\section*{一、开场（0—5 分钟）}

明说：今天不学新东西，只把 14 节课串成一条线。

一句话铺垫：所有内容你都见过，只是散着放；今天给它们找一个共同的名字。

\section*{二、板书主线图（5—26 分钟）}

\begin{caozuo}
  在黑板正中间写一个大字「\textbf{群}」，往三个方向拉线。三条线一边写一边讲，
  每写一个点就回头问一句「这个我们什么时候见过」。
\end{caozuo}

\textbf{左边一条线：数着数着，就绕回来了。}

\begin{itemize}
  \item 模 \(12\) 加法：\(8+7=3\)，走 \(7\) 格，绕回来了；
  \item 单位元 \(0\)：加上它等于没加；
  \item 逆元：\(1\) 与 \(11\)、\(2\) 与 \(10\)、\(6\) 与自己；
  \item 由 \(3\) 生成的 \(\{0,3,6,9\}\)：只走 \(4\) 步就回到 \(0\)。
\end{itemize}

\textbf{右边一条线：位置换一换，还是那几个样子。}

\begin{itemize}
  \item 正三角形：\(6\) 个对称动作，记作 \(S_3\)；
  \item 正方形：\(8\) 个对称动作，记作 \(D_4\)；
  \item 复合当作运算：先 \(r\) 后 \(f\) 与先 \(f\) 后 \(r\) 不一样，\(rf\neq fr\)；
  \item 阶与子群：子群的大小整除群的大小。
\end{itemize}

\textbf{下方一条线：幂的尾巴，也绕圈。}

\begin{itemize}
  \item \(2^n\) 除以 \(7\)，余数在 \(2,4,1\) 里打转；
  \item 由现象归纳出猜想，再证明它：费马小定理；
  \item 回头回答第 1 课：\(10^{6}\equiv1\pmod 7\)，所以 \textbf{\(7\) 整除 \(999999\)}。
\end{itemize}

\begin{caozuo}
  三条线写完，在「群」旁边补上四条规则，让全班一起念一遍。
  落一句：数、动作、置换，只要满足这四条，数学家就把它们归到一起研究。
\end{caozuo}

\begin{definition}[群]
  一个集合配上一种运算，如果满足下面四条，就叫\textbf{群}：
  \begin{enumerate}
    \item \textbf{封闭}：任取 \(a,b\)，\(a*b\) 还在这里面；
    \item \textbf{结合}：\((a*b)*c=a*(b*c)\)；
    \item \textbf{单位元}：存在 \(e\)，使 \(a*e=e*a=a\)；
    \item \textbf{逆元}：每个 \(a\) 都有 \(b\)，使 \(a*b=b*a=e\)。
  \end{enumerate}
\end{definition}

\section*{三、抽查（26—36 分钟）}

\begin{caozuo}
  随机点座位号，指着主线图上的一个点问「这是什么意思？举个例子」。
  点 8—10 个，答不上来就往前回溯一句，不追究、不批评。
\end{caozuo}

\begin{zhu}
  抽到的题尽量覆盖三条线，各来两三道。用手指回答也算数。四道抽查题的答案：
  \begin{enumerate}
    \item 模 \(12\) 加法里 \(5\) 的逆元是 \textbf{\(7\)}（\(5+7=12\equiv0\)）；
    \item 正三角形有 \(6\) 个对称动作，正方形有 \textbf{\(8\)} 个；
    \item 三角形的对称里 \(rf\) 与 \(fr\) \textbf{不一样}（\(rf\neq fr\)）；
    \item \(2^{100}\) 除以 \(7\) 余 \textbf{\(2\)}（\(100=6\times16+4\)，所以
      \((2^{6})^{16}\times2^{4}\equiv1\times16\equiv2\)）。
  \end{enumerate}
\end{zhu}

\section*{四、展望（36—45 分钟）}

\begin{caozuo}
  在黑板角落写「加法 \(+\) 乘法，放一起」。\textbf{先回指一句分配律}（就是下面那句），
  学生早就在用了，不必从零解释；再给名字和一个例子，不展开。
\end{caozuo}

\noindent 分配律你们早就在用了：\(3\times(4+5)=3\times4+3\times5\)。

\begin{definition}[环]
  一个集合上同时有加法和乘法，如果下面三条都满足，就叫做\textbf{环}：
  \begin{enumerate}
    \item \textbf{加法}：满足第 4 课那四条规则（封闭、结合、单位元、逆元），而且加法可以交换：\(a+b=b+a\)；
    \item \textbf{乘法}：封闭、结合、\textbf{有单位元 \(1\)}（\(a\cdot1=1\cdot a=a\)）、而且乘法可以交换：\(a\cdot b=b\cdot a\)——但乘法不要求每个数都有逆元；
    \item \textbf{乘法对加法有分配律}：\(a\cdot(b+c)=a\cdot b+a\cdot c\)，\((b+c)\cdot a=b\cdot a+c\cdot a\)。
  \end{enumerate}
  乘法可以交换，所以它又叫\textbf{交换环}——乘法不交换的环也有，这门课不去碰它。
\end{definition}

\begin{definition}[域]
  在上面这个环的基础上，如果每个非零元素都有乘法逆元（能做除法），就叫做\textbf{域}。
\end{definition}

\begin{example}
  整数 \(\mathbb{Z}\) 是环但不是域（\(2\) 没有整数逆元，\(\tfrac12\notin\mathbb{Z}\)）；
  有理数 \(\mathbb{Q}\) 是域。
\end{example}

\begin{caozuo}
  \textbf{（附）三次、四次方程长什么样——每页一个，共三页，只给样子。}
  三页的顺序：一元二次（配方 \(+\) 判别式）\(\to\) 一元三次（卡尔达诺设 \(t=u+v\)）
  \(\to\) 一元四次（费拉里配两次方，中间变量 \(y\)）。三页都写成
  「先设中间变量」的样子，\textbf{不要展开成一长条完整求根公式}。
\end{caozuo}

\begin{example}[一元二次方程：配方就能解]
  \[ ax^2+bx+c=0\quad(a\neq0) \]
  两边除以 \(a\)，再令中间变量
  \[ t=x+\frac{b}{2a}, \]
  就化成
  \[ t^2=\frac{b^2-4ac}{4a^2}. \]
  令 \(\Delta=b^2-4ac\)（叫\textbf{判别式}），于是
  \[ t=\pm\frac{\sqrt{\Delta}}{2a},\qquad x=\frac{-b\pm\sqrt{b^2-4ac}}{2a}. \]
  \textbf{一句话}：先把 \(x\) 平移到 \(t\)，方程就只剩一个平方。
  （\(\Delta>0\) 两个根，\(\Delta=0\) 一个根，\(\Delta<0\) 没有实数根。）
\end{example}

\begin{example}[一元三次方程：卡尔达诺设 \(t=u+v\)]
  \[ t^3+pt+q=0 \]
  （任何三次方程都能先平移消掉平方项，变成这个「缺项」的样子。）
  令
  \[ t=u+v,\qquad\text{并约定}\ uv=-\frac{p}{3}, \]
  代入后方程变成
  \[ u^3+v^3=-q, \]
  而由约定又有
  \[ u^3v^3=-\frac{p^3}{27}. \]
  所以 \(u^3\) 与 \(v^3\) 是一元二次方程
  \[ z^2+qz-\frac{p^3}{27}=0 \]
  的两个根。\textbf{解出这个二次方程，就得到 \(u^3,v^3\)，再开立方就得到 \(t=u+v\)。}

  \textbf{一句话}：三次方程被「两个中间变量」降成了一个二次方程。
\end{example}

\begin{example}[一元四次方程：费拉里配两次方]
  \[ t^4+pt^2+qt+r=0 \]
  （同样先平移消掉三次项。）第一步配方：
  \[ \left(t^2+\frac{p}{2}\right)^2=-\frac{p^2}{4}+r-qt. \]
  第二步，\textbf{引入中间变量 \(y\)}，给左边添上 \(2y\left(t^2+\frac p2\right)+y^2\)，
  让它仍然是完全平方：
  \[ \left(t^2+\frac{p}{2}+y\right)^2
     =2y\,t^2-qt+\left(y^2+py-\frac{p^2}{4}+r\right). \]
  只要右边的判别式为 \(0\)（即 \(q^2=8y^3+8py^2-8ry\)，这是\textbf{一个关于 \(y\) 的三次方程}），
  右边也是完全平方；于是两边开方，就得到\textbf{两个二次方程}。

  \textbf{一句话}：四次方程靠「中间变量 \(y\)」拆成两个二次方程——而 \(y\) 又要靠上面那个三次方程求出来。
\end{example}

\begin{zhu}
  这 3 页\textbf{只要求「看个样子」，不要求记、不要求算}。讲的时候一句话带过：
  「二次你们会，三次四次十六世纪就有人做出来了，长这样——注意它有多长。」
  \textbf{建议总用时不超过 4 分钟}；若时间紧，只讲二次那一页，三次四次两页让学生课后看。
  学生版这三页各留了一道横线，是给他们随手写两个记号用的，不必收上来。
  \(\sqrt{\ }\) 是平方根、\(\sqrt[3]{\ }\) 是立方根，学生见过但没系统学过，
  只要说「认得这个记号就行」，不要展开讲开方。
\end{zhu}

\begin{caozuo}
  讲一个故事：一元二次方程有求根公式，三次、四次也有。数学家找了三百多年，
  一直到十九世纪才有人证明：五次方程根本没有这样的公式。做这件事的人叫伽罗瓦，
  二十岁就去世了。他用到的工具，就是这门课讲的群。
\end{caozuo}

\begin{xiaojie}
  \begin{itemize}
    \item 这门课看了三类「绕圈」：数的加法、图形的对称、幂的余数。
    \item 它们满足同样的四条规则，所以共用同一个名字——群。
    \item 群之后还有环与域；五次方程的故事，是群给出的答案。
  \end{itemize}
\end{xiaojie}

\noindent\textbf{黑板上的最后一行}：群不是新东西，它就是「绕圈」这件事的名字。而它，解开了三百年的谜。